% Part: normal-modal-logic
% Chapter: filtrations
% Section: more-filtrations

\documentclass[../../../include/open-logic-section]{subfiles}

\begin{document}

\olfileid[pt]{nml}{fil}{acc}

\olsection{Filtrações e propriedades de acessibilidade}

Como observamos, filtrações de modelos universais, seriais e reflexivos
são sempre também universais, seriais ou reflexivas. Mas nem toda
filtração de um modelo simétrico ou transitivo é, respectivamente,
simétrica ou transitiva. Em alguns casos, porém, é possível definir
filtrações de modo que isso valha. Para fazê-lo, procedemos como na
definição da filtração mais grossa, mas acrescentamos condições à
definição de~$R^*$. Seja $\Gamma$ fechado sob sub-!!{formula}s. Considere
as relações~$C_i(u,v)$ de \olref{tab:Cn-filtrations} entre mundos $u$,
$v$ em um modelo $\mModel{M} =\tuple{W, R, V}$. Podemos definir
$R^*[u][v]$ com base em combinações dessas condições. Por exemplo, se
estipulamos que $R^*[u][v]$ sse vale a condição $C_1(u,v)$, obtemos
exatamente a filtração mais grossa. Se estipulamos $R^*[u][v]$ sse valem
tanto $C_1(u,v)$ quanto $C_2(u, v)$, obtemos uma filtração diferente.
Ela é ``mais fina'' que a mais grossa, pois menos pares de mundos
satisfazem $C_1(u,v)$ e $C_2(u,v)$ do que apenas $C_1(u,v)$.

\begin{table}[ht]
  \centering
  \begin{tabular}{|ll|}
    \hline
    \multirow{2}{*}{$C_1(u,v)$:}
    \iftag{prvBox}{&
      se $\Box!A \in \Gamma$ e $\mSat{M}{\Box!A}[u]$,
      então $\mSat{M}{!A}[v]$;\iftag{prvDiamond}{ e}{}\\}{}
    \iftag{prvDiamond}{&
      se $\Diamond!A \in \Gamma$ e $\mSat{M}{!A}[v]$,
      então $\mSat{M}{\Diamond!A}[u]$; \\}{}
    \hline
    \multirow{2}{*}{$C_2(u,v)$:} 
    \iftag{prvBox}{&
      se $\Box!A \in \Gamma$ e $\mSat{M}{\Box!A}[v]$,
      então $\mSat{M}{!A}[u]$;\iftag{prvDiamond}{ e}{}\\}{}
    \iftag{prvDiamond}{&
      se $\Diamond!A \in \Gamma$ e $\mSat{M}{!A}[u]$,
      então $\mSat{M}{\Diamond!A}[v]$; \\}{}
    \hline
    \multirow{2}{*}{$C_3(u,v)$:}
    \iftag{prvBox}{&
      se $\Box!A \in \Gamma$ e $\mSat{M}{\Box!A}[u]$,
      então $\mSat{M}{\Box!A}[v]$;\iftag{prvDiamond}{ e}{}\\}{}
    \iftag{prvDiamond}{&
      se $\Diamond!A \in \Gamma$ e $\mSat{M}{\Diamond!A}[v]$,
      então $\mSat{M}{\Diamond!A}[u]$; \\}{}
    \hline
    \multirow{2}{*}{$C_4(u,v)$:}
    \iftag{prvBox}{&
      se $\Box!A \in \Gamma$ e $\mSat{M}{\Box!A}[v]$,
      então $\mSat{M}{\Box!A}[u]$;\iftag{prvDiamond}{ e}{}\\}{}
    \iftag{prvDiamond}{&
      se $\Diamond!A \in \Gamma$ e $\mSat{M}{\Diamond!A}[u]$,
      então $\mSat{M}{\Diamond!A}[v]$; \\}{}
    \hline
    \end{tabular}
  \caption{Condições sobre mundos possíveis para definir
    filtrações.}
  \ollabel{tab:Cn-filtrations}
\end{table}

\begin{thm}\ollabel{thm:more-filtrations}
  Seja $\mModel{M} =\tuple{W,R,V}$ um modelo e seja $\Gamma$ fechado sob
  sub-!!{formula}s. Sejam $W^*$ e $V^*$ definidos como em
  \olref[fil]{defn:filtration}. Então:
  \begin{enumerate}
  \item Suponha $R^*[u][v]$ se, e somente se, $C_1(u, v) \land
    C_2(u,v)$. Então $R^*$ é simétrica e $\mModel{M^*} =
    \tuple{W^*,R^*,V^*}$ é uma filtração se $\mModel{M}$ é simétrico.
  \item Suponha $R^*[u][v]$ se, e somente se, $C_1(u, v)
    \land C_3(u,v)$. Então $R^*$ é transitiva e
    $\mModel{M^*}=\tuple{W^*,R^*,V^*}$ é uma filtração se $\mModel{M}$
    é transitivo.
  \item Suponha $R^*[u][v]$ se, e somente se, $C_1(u, v) \land C_2(u,v)
    \land C_3(u,v) \land C_4(u,v)$. Então $R^*$ é simétrica e
    transitiva, e $\mModel{M^*}=\tuple{W^*,R^*,V^*}$ é uma filtração
    se $\mModel{M}$ é simétrico e transitivo.
  \item Suponha que $R^*$ seja definida por $R^*[u][v]$ se, e somente
    se, $C_1(u, v) \land C_3(u,v) \land C_4(u,v)$. Então $R^*$ é
    transitiva e euclidiana, e $\mModel{M^*}=\tuple{W^*,R^*,V^*}$ é uma
    filtração se $\mModel{M}$ é transitivo e euclidiano.
  \end{enumerate}
\end{thm}

\begin{proof}
  \begin{enumerate}
    \item É imediato que $R^*$ é simétrica, pois $C_1(u,v)
      \Leftrightarrow C_2(v,u)$ e $C_2(u,v) \Leftrightarrow
      C_1(v,u)$. Resta mostrar que, se $\mModel{M}$ é simétrico, então
      $\mModel{M^*}$ é uma filtração por meio de~$\Gamma$. A condição
      $C_1(u,v)$ garante o cumprimento de
      \iftag{prvBox}{%
        \iftag{prvDiamond}{%
          \olref[fil]{defn:filtration-R2} e
          \olref[fil]{defn:filtration-R3} de
          \olref[fil]{defn:filtration}}{%
          \olref[fil]{defn:filtration}\olref[fil]{defn:filtration-R2}}}{%
        \olref[fil]{defn:filtration}\olref[fil]{defn:filtration-R3}}.
      Assim, precisamos apenas verificar
      \olref[fil]{defn:filtration}\olref[fil]{defn:filtration-R1},
      O.s.\@, que $Ruv$ implica $R^*[u][v]$.

      Suponha, portanto, $Ruv$. Para mostrar $R^*[u][v]$, precisamos
      estabelecer $C_1(u,v)$ e $C_2(u,v)$. Para $C_1$:
      \iftag{prvBox}{se $\Box!A \in\Gamma$ e $\mSat{M}{\Box!A}[u]$,
        então também $\mSat{M}{!A}[v]$ (pois $Ruv$).
        \iftag{prvDiamond}{ De modo semelhante, }{}}{}
      \iftag{prvDiamond}{se $\Diamond!A \in \Gamma$ e
        $\mSat{M}{!A}[v]$, então $\mSat{M}{\Diamond!A}[u]$, pois
        $Ruv$.}{} Para $C_2$: \iftag{prvBox}{se $\Box!A \in \Gamma$ e
        $\mSat{M}{\Box!A}[v]$, então $Ruv$ implica $Rvu$ por simetria;
        logo, $\mSat{M}{!A}[u]$.\iftag{prvDiamond}{ De modo semelhante,
        }{}}{}\iftag{prvDiamond}{se $\Diamond!A \in\Gamma$ e
        $\mSat{M}{!A}[u]$, então $\mSat{M}{\Diamond!A}[v]$ (pois $Rvu$
        por simetria).}{}
    \item Exercício.
    \item Exercício.
    \item Exercício.
  \end{enumerate}
\end{proof}

\begin{prob}
  Complete a demonstração de \olref[nml][fil][acc]{thm:more-filtrations}.
\end{prob}

\end{document}
